\section*{अध्याय \ref{ch:g9:natural-risks} --- प्राकृतिक जोखिमों के साथ जीना}

\begin{solution}{exo:g9:natural-risks:1}
(क) संकट। (ख) \omterm{def:g9:natural-risks:vulnerability}{सुभेद्यता}। (ग) \omterm{def:g9:natural-risks:vulnerability}{सुभेद्यता}। (घ) संकट। (ङ) दोनों: कीचड़ प्रवाह संकट हैं;
उनके रास्ते में कस्बा बसाना उसे सुभेद्य बनाता
है।
\end{solution}

\begin{solution}{exo:g9:natural-risks:2}
(क) ग़लत: जहाँ कोई नहीं रहता, वहाँ किसी चीज़ को नुक़सान नहीं पहुँच सकता, इसलिए \omterm{def:g9:natural-risks:risk}{जोखिम} कम है।
(ख) सही। (ग) सही। (घ) ग़लत: चेतावनी \omterm{def:g4:volcanoes-earthquakes-first-look:volcano}{ज्वालामुखी} को नहीं बदलती; वह लोगों की \omterm{def:g9:natural-risks:vulnerability}{सुभेद्यता}
घटाती है, जो समय पर निकल सकते हैं।
\end{solution}

\begin{solution}{exo:g9:natural-risks:3}
(क) हाँ, दिनों पहले। (ख) हाँ, घंटों से दिनों पहले। (ग) नहीं: \omterm{def:g4:volcanoes-earthquakes-first-look:earthquake}{भूकंपों} की भविष्यवाणी नहीं
की जा सकती। (घ) हाँ: \omterm{def:g7:earthquakes:tsunami}{सुनामी} शुरू होने के बाद \omterm{def:g2:water-on-earth:ocean}{महासागर} पार करने में उसे घंटों लगते हैं, और
चेतावनी उससे तेज़ चलती है।
\end{solution}

\begin{solution}{exo:g9:natural-risks:4}
(क) ज़मीन का हिलना: \omterm{def:g7:volcanoes:magma}{मैग्मा} के ऊपर उठने पर छोटे \omterm{def:g4:volcanoes-earthquakes-first-look:earthquake}{भूकंपों} के झुंड। (ख) ढलान का झुकाव:
\omterm{def:g4:volcanoes-earthquakes-first-look:volcano}{ज्वालामुखी} का फूलना। (ग) ज़मीन मिलीमीटरों में कैसे खिसकती है: \omterm{def:g7:volcanoes:magma}{मैग्मा} के इकट्ठा होने पर
फूलना। (घ) धुएँ के गुबार की गैसें: \omterm{def:g7:volcanoes:magma}{मैग्मा} के ऊपर आने पर ज़्यादा गैस का
निकलना।
\end{solution}

\begin{solution}{exo:g9:natural-risks:5}
(क) $45 \div 2 = \qty{22.5}{km/h}$। (ख) $4 \times 2 = \qty{8}{km}$, ऊँची ज़मीन तक
पहुँचने के लिए काफ़ी। (ग) \omterm{def:g5:natural-hazards:storm}{तूफ़ान} में छिपे \omterm{def:g4:volcanoes-earthquakes-first-look:eruption}{उद्गार} को किसी ने नहीं देखा, और किसी ने उन्हें
चेतावनी नहीं दी: कस्बे तक जानकारी जल्दी भेजने का कोई तरीक़ा नहीं था, और कोई निकासी योजना
नहीं थी। (घ) उदाहरण के लिए: ऊँची ज़मीन तक के रास्तों वाली एक निकासी योजना, जिसका पहले से
अभ्यास हुआ हो; वैज्ञानिकों से कस्बे तक चेतावनी की एक कड़ी; \omterm{def:g4:volcanoes-earthquakes-first-look:volcano}{ज्वालामुखी} पर ज़्यादा
उपकरण।
\end{solution}

\begin{solution}{exo:g9:natural-risks:6}
(क) \omterm{def:g1:landscapes-around-us:water}{नदी} से \qty{10}{m} ऊपर की \omterm{def:g5:water-shapes-the-land:flood}{बाढ़}। (ख) \qty{2}{m} की \omterm{def:g5:water-shapes-the-land:flood}{बाढ़} के लिए लगभग
\qty{0.47}{km}, \qty{5}{m} की \omterm{def:g5:water-shapes-the-land:flood}{बाढ़} के लिए \qty{0.56}{km} और
\qty{10}{m} की \omterm{def:g5:water-shapes-the-land:flood}{बाढ़} के लिए \qty{0.67}{km} (लगभग
\qty{0.05}{km} के भीतर के पाठ्यांक ठीक हैं)। (ग) केवल \qty{10}{m} की \omterm{def:g5:water-shapes-the-land:flood}{बाढ़}: ज़मीन
\omterm{def:g1:landscapes-around-us:water}{नदी} से 2 और \qty{5}{m} से ज़्यादा ऊपर है। (घ) उस तक पहुँचने के लिए \omterm{def:g1:landscapes-around-us:water}{नदी} का थोड़ा-सा उठना
काफ़ी है, और छोटी \omterm{def:g5:water-shapes-the-land:flood}{बाढ़ें} बड़ी \omterm{def:g5:water-shapes-the-land:flood}{बाढ़ों} से ज़्यादा बार आती
हैं।
\end{solution}

\begin{solution}{exo:g9:natural-risks:7}
(क) ऊँचा संकट, कम \omterm{def:g9:natural-risks:vulnerability}{सुभेद्यता}: मध्यम \omterm{def:g9:natural-risks:risk}{जोखिम}। (ख) कम संकट, ऊँची \omterm{def:g9:natural-risks:vulnerability}{सुभेद्यता}: मध्यम \omterm{def:g9:natural-risks:risk}{जोखिम}।
(ग) ऊँचा संकट, ऊँची \omterm{def:g9:natural-risks:vulnerability}{सुभेद्यता}: ऊँचा \omterm{def:g9:natural-risks:risk}{जोखिम}। (घ) कम संकट, कम \omterm{def:g9:natural-risks:vulnerability}{सुभेद्यता}: कम
\omterm{def:g9:natural-risks:risk}{जोखिम}।
\end{solution}

\begin{solution}{exo:g9:natural-risks:8}
(क) C, फिर B, फिर A। (ख) टेक त्रिभुज बनाते हैं, जो किसी टेक को खींचे या दबाए बिना झुक नहीं
सकते: ढाँचा झुकने के बजाय अपना आकार बनाए रखता है। (ग) अपने नीचे ज़मीन के आगे-पीछे खिसकने
पर विलगक लचक जाते हैं, इसलिए गति का बहुत कम हिस्सा स्कूल तक पहुँचता
है।
\end{solution}

\begin{solution}{exo:g9:natural-risks:9}
(क) $1 \div 100 = 1~\%$ प्रति वर्ष। (ख) संभावना हर वर्ष वही रहती है, पिछले वर्ष चाहे जो
हुआ हो: \omterm{def:g5:water-shapes-the-land:flood}{बाढ़} अगले वर्ष फिर आ सकती है।
(ग) $300 \div 100 = 3$ \omterm{def:g5:water-shapes-the-land:flood}{बाढ़ें}, औसतन।
\end{solution}

\begin{solution}{exo:g9:natural-risks:10}
(क) $3600 \div 15 = 240$ माप। (ख) $240 \times 24 \times 365 =
\num{2102400} \approx 2.1 \times 10^6$ माप। (ग) $2400 \div 800 =
\qty{3}{h}$। (घ) इन तीन घंटों में चेतावनी केंद्र \omterm{def:g1:landscapes-around-us:water}{तट} को
चेतावनी दे सकते हैं और लोग ऊँची ज़मीन पर जा सकते हैं।
\end{solution}

\begin{solution}{exo:g9:natural-risks:11}
(क) $4 \times 3 \times 3.8 = \qty{45.6}{L}$। (ख) $45.6 \div 1.5 = 30.4$,
इसलिए 31 बोतलें। (ग) उदाहरण के लिए: टिकने वाला भोजन, दुकानें बंद हों तब खाने के लिए; एक
टॉर्च, अगर बिजली कट जाए; बैटरी या हाथ से चलने वाला रेडियो, चेतावनियाँ सुनने के लिए;
प्राथमिक चिकित्सा किट, छोटी चोटों के इलाज के लिए; गरम कंबल। (घ) ताकि अगर परिवार को जल्दी
निकलना पड़े तो उसे तुरंत ले जाया जा सके।
\end{solution}

\begin{solution}{exo:g9:natural-risks:12}
उदाहरण के लिए: (1) कीचड़ प्रवाहों का एक \omterm{def:g9:natural-risks:hazard-map}{संकट मानचित्र} बनाना, घाटी-दर-घाटी
(संकट वाली ओर); (2) \omterm{def:g4:volcanoes-earthquakes-first-look:volcano}{ज्वालामुखी} पर \omterm{def:g7:earthquakes:seismometer}{भूकंपमापी}, झुकावमापी, GPS अभिग्राही और गैस
उपकरण लगाना (संकट वाली ओर); (3) वेधशाला से कस्बे तक, सायरनों के साथ, चेतावनी की एक
तेज़ कड़ी बनाना (\omterm{def:g9:natural-risks:vulnerability}{सुभेद्यता} वाली
ओर); (4) ऊँची ज़मीन तक के रास्तों की योजना बनाना, उन पर निशान लगाना, और अभ्यासों से उनका
अभ्यास करना (\omterm{def:g9:natural-risks:vulnerability}{सुभेद्यता} वाली ओर); (5) कीचड़ प्रवाहों के रास्ते में कोई नया घर या स्कूल न
बनाना (\omterm{def:g9:natural-risks:vulnerability}{सुभेद्यता} वाली ओर)।
\end{solution}

\begin{solution}{pb:g9:natural-risks:1}
\textbf{1.} \omterm{def:g1:landscapes-around-us:water}{नदी} C के बगल में \qty{143}{m} पर, A के बगल में \qty{137}{m} पर
और B के बगल में \qty{126}{m} पर है; \omterm{def:g2:drawing-to-map:map}{मानचित्र} पर C पश्चिम में है और B पूर्व में,
इसलिए \omterm{def:g1:landscapes-around-us:water}{नदी} पूर्व की ओर बहती है।

\textbf{2.} A: $137 - 137 = \qty{0}{m}$; B: $276 - 126 = \qty{150}{m}$;
C: $236 - 143 = \qty{93}{m}$।

\textbf{3.} केवल A, जो \omterm{def:g1:landscapes-around-us:water}{नदी} से \qty{5}{m} से कम ऊपर है।

\textbf{4.} लगभग \qty{0.57}{km} (0.5 और \qty{0.65}{km} के बीच)।

\textbf{5.} \qty{100}{m} पर उठाव \qty{200}{m} पर उठाव का आधा
है: A $1.4 \div 2 = 0.7~\%$; B $83 \div 2 = 41.5~\%$;
C $20 \div 2 = 10~\%$।

\textbf{6.} B: उसकी ढलान बहुत खड़ी है, और ढीली \omterm{def:g3:stones-pebbles-sand:rock}{चट्टान} और मिट्टी खड़ी ढलान पर नीचे खिसक
सकती है, ख़ासकर भारी \omterm{def:g4:clouds-rain-snow:precipitation}{बारिश} के बाद।

\textbf{7.} A: लगभग \qty{1}{km}; B: लगभग \omterm{def:g7:earthquakes:fault}{भ्रंश} पर (\qty{0.05}{km} से
कम); C: लगभग \qty{1.3}{km}।

\textbf{8.} अगर \omterm{def:g7:earthquakes:fault}{भ्रंश} खिसकता है, तो ज़मीन उसके साथ-साथ टूटती और सरकती है, जो किसी इमारत को
चीर सकती है; और \omterm{def:g7:earthquakes:fault}{भ्रंश} पर आया \omterm{def:g4:volcanoes-earthquakes-first-look:earthquake}{भूकंप} वहाँ की ज़मीन को बहुत ज़ोर से
हिलाएगा।

\textbf{9.}

\begin{center}
{\small
\begin{tabular}{|c|c|c|c|}
\hline
जगह & \omterm{def:g5:water-shapes-the-land:flood}{बाढ़} क्षेत्र में? & खड़ी ढलान? & \omterm{def:g7:earthquakes:fault}{भ्रंश} पर? \\
\hline
A & हाँ & नहीं & नहीं \\
B & नहीं & हाँ & हाँ \\
C & नहीं & नहीं & नहीं \\
\hline
\end{tabular}}
\end{center}

\textbf{10.} जगह C: वही एक जगह है जिसके तीनों जवाब ``नहीं'' हैं: \omterm{def:g1:landscapes-around-us:water}{नदी} से \qty{93}{m}
ऊपर, \qty{10}{\%} की हल्की ढलान, \omterm{def:g7:earthquakes:fault}{भ्रंश} से \qty{1}{km} से ज़्यादा
दूर।

\textbf{11.} उदाहरण के लिए: ढाँचे में तिरछे टेक और अपरूपण दीवारें, मज़बूत जोड़, या इमारत
के नीचे विलगक; छात्र \omterm{def:g4:volcanoes-earthquakes-first-look:earthquake}{भूकंप} के अभ्यास करते हैं (झुको, ढको, पकड़े
रहो)।

\textbf{12.} 28 में से \textbf{17 घर} \omterm{def:g5:water-shapes-the-land:flood}{बाढ़} क्षेत्र के अंदर खड़े हैं:
$17 \div 28 \approx 0.61$, लगभग 5 में से 3 घर।
\end{solution}
